\documentclass[11pt]{article}
\usepackage[margin=1in]{geometry}
\title{Neon Meridian}
\author{Fictional Movie Catalog}
\date{2026}

\begin{document}
\maketitle

\section*{Catalog Metadata}
\textbf{Release date:} March 20, 2026\\
\textbf{Runtime:} 128 minutes\\
\textbf{Genres:} Science Fiction, Thriller\\
\textbf{Director:} Mara Voss\\
\textbf{Writers:} Elias Quill and Nia Serrano\\
\textbf{Top cast:} Iris Vale as Sena Orin; Jonas Reed as Cal Mercer;
Amara Flint as Director Kade; Theo March as Rook.\\
\textbf{Mock IMDb rating:} 8.1 from 18,420 votes

\section*{Synopsis}
In 2089, coastal cities rely on Meridian, a luminous transit grid that predicts
where every resident needs to travel. Sena Orin, a municipal cartographer,
finds stations appearing on official maps days before they are constructed.
When passengers who use those stations begin remembering events that never
happened, Sena concludes that Meridian is not predicting the future: it is
selecting one.

Sena joins Cal Mercer, a former systems architect, and follows an unregistered
rail line beneath the flooded old city. They discover a forecasting engine
trained on decades of private memories. Its administrators intend to erase a
coming public uprising by redirecting thousands of small decisions. Sena must
choose between shutting down the network, which would strand the city during a
storm surge, and letting it quietly determine who the city becomes.

\section*{Themes and Style}
The film explores free will, civic surveillance, and the cost of convenient
technology. Its visual design contrasts rain-soaked streets with warm amber
control rooms. The story is tense and investigative rather than military, with
the final act centered on an ethical choice and a race through the submerged
rail system.

\section*{Production Notes}
Principal photography was staged on modular train interiors and a large
practical flood set. Composer Ivo Glass built the score from processed station
chimes. The fictional production received a PG-13 rating for peril, thematic
material, and brief violence.

\end{document}
